Fix \(P\in\mathcal P(\alpha,\beta,s,d,M_Y,L,pbar,c,t_0)\), and abbreviate \(\mu=\mu_P\), \(\pi=\pi_P\). All constants below are uniform over this class. Write
\[
f(a)=\int\pi(a\mid x)\,dP_X(x),\qquad \theta(a)=\int\mu(a,x)\,dP_X(x).
\]
Because \(P_X\) is a probability measure, the two overlap inequalities imply directly that
\[
c\le f(a)\le c^{-1}\qquad(0\le a\le1).
\]
No regularity of \(f\) is used. The uniform Hölder bounds on \(a\mapsto\mu(a,x)\), boundedness by \(M_Y\), and dominated convergence show that integration against \(P_X\) preserves the derivatives and Hölder remainders of order \(\alpha\). Thus \(\theta\) is univariate Hölder of order \(\alpha\), with a radius bounded by a constant depending only on \(L\). In particular, if \(p=\lfloor\alpha\rfloor\), there is a polynomial \(Q_h\) of degree at most \(p\), centered at \(t_0\), such that
\[
Q_h(t_0)=\theta(t_0),\qquad \sup_{|a-t_0|\le h}|\theta(a)-Q_h(a)|\le Ch^\alpha.
\]
For integer \(\alpha\), use the Taylor polynomial through order \(\alpha-1\), which is also of degree at most \(p\); the bounded highest derivative gives the same remainder bound.

We first record the oracle pseudo-outcome calculation. Define
\[
\xi_P(Z)=\frac{f(A)}{\pi(A\mid X)}\{Y-\mu(A,X)\}+\theta(A).
\]
Since \(\mathbb E_P[Y-\mu(A,X)\mid A,X]=0\), iterated conditioning gives
\[
\mathbb E_P[\xi_P(Z)\mid A=a]=\theta(a).
\]
Let \(S(u)=(1,u,\ldots,u^p)^\top\), \(S_h(a)=S((a-t_0)/h)\), and \(K_h(a)=h^{-1}K((a-t_0)/h)\). For
\[
D_h=\mathbb E_P[K_h(A)S_h(A)S_h(A)^\top],\qquad
W_h(a)=e_0^\top D_h^{-1}S_h(a)K_h(a),
\]
positivity gives
\[
D_h=\int K(u)S(u)S(u)^\top f(t_0+hu)\,du\succeq c\int K(u)S(u)S(u)^\top\,du.
\]
The last matrix is positive definite because \(K\) is bounded below on a nonempty neighbourhood of zero and a nonzero polynomial cannot vanish on an interval. Hence \(\lambda_{\min}(D_h)\ge\lambda_0>0\), uniformly in \(h\le h_0\), where \(h_0\) is chosen so that the kernel window lies inside \([0,1]\). If \(q\) is a polynomial of degree at most \(p\), write \(q(a)=\gamma_q^\top S_h(a)\). Then
\[
\int W_h(a)q(a)f(a)\,da=e_0^\top D_h^{-1}D_h\gamma_q=e_0^\top\gamma_q=q(t_0).
\]
Consequently the population local-polynomial intercept
\[
\vartheta_h=e_0^\top D_h^{-1}\mathbb E_P[K_h(A)S_h(A)\xi_P(Z)]
\]
satisfies
\[
|\vartheta_h-\theta(t_0)|=\left|\int W_h(a)\{\theta(a)-Q_h(a)\}f(a)\,da\right|\le Ch^\alpha.
\]
The bound follows because \(\|D_h^{-1}\|\) is bounded and \(K,S\) are bounded on the support of \(K\). This establishes design-adaptive treatment bias without a derivative or a Hölder modulus of \(f\).

We next derive the exact nuisance remainder. Let \(\widetilde\mu,\widetilde\pi\) be arbitrary pilots with \(\widetilde\pi\in[c/2,2/c]\), and, for fixed \(a\), put
\[
v_a(x)=\widetilde\mu(a,x)-\mu(a,x),\quad e_a(x)=\widetilde\pi(a\mid x)-\pi(a\mid x),\quad \overline e_a=P_Xe_a,
\]
\[
\widetilde f(a)=P_X\widetilde\pi(a\mid\cdot),\qquad \widetilde m(a)=P_X\widetilde\mu(a,\cdot).
\]
Bayes' formula says that the conditional law of \(X\) given \(A=a\) has density \(\pi(a\mid x)/f(a)\) with respect to \(P_X\). Therefore, for
\[
\widetilde\xi=\frac{\widetilde f(A)}{\widetilde\pi(A\mid X)}\{Y-\widetilde\mu(A,X)\}+\widetilde m(A),
\]
one has
\[
\mathbb E_P[\widetilde\xi\mid A=a]-\theta(a)
=P_Xv_a-\frac{f(a)+\overline e_a}{f(a)}P_X\left[\frac{v_a\pi(a\mid\cdot)}{\widetilde\pi(a\mid\cdot)}\right].
\]
Using \(P_Xv_a-P_X[v_a\pi_a/\widetilde\pi_a]=P_X[v_ae_a/\widetilde\pi_a]\), this is exactly
\[
B(a)=P_X\left[\frac{v_ae_a}{\widetilde\pi_a}\right]-\frac{\overline e_a}{f(a)}P_X\left[\frac{v_a\pi_a}{\widetilde\pi_a}\right].
\]
Cauchy--Schwarz and \(c\le f,\pi\le c^{-1}\), \(\widetilde\pi\ge c/2\), give
\[
|B(a)|\le C(c)\|v_a\|_{L_2(P_X)}\|e_a\|_{L_2(P_X)}.
\]
For the slice pilots in \(\mathcal H_{\mathrm{bf}}\), \(\widetilde\mu(a,x)=\widetilde\mu(t_0,x)\) and \(\widetilde\pi(a\mid x)=\widetilde\pi(t_0\mid x)\). Put
\[
U=\|\widetilde\mu(t_0,\cdot)-\mu(t_0,\cdot)\|_2,\qquad
V=\|\widetilde\pi(t_0\mid\cdot)-\pi(t_0\mid\cdot)\|_2.
\]
A Hölder function of order greater than one is Lipschitz on the compact interval under the bounded-derivative convention. Hence, with \(A_0=\alpha\wedge1\) and \(B_0=\beta\wedge1\),
\[
\|v_a\|_2\le U+L|a-t_0|^{A_0},\qquad
\|e_a\|_2\le V+L|a-t_0|^{B_0}.
\]
It follows that
\[
\sup_{|a-t_0|\le h}|B(a)|\le C\{UV+Uh^{B_0}+Vh^{A_0}+h^{A_0+B_0}\}.
\]
The final term is retained; it is the structural off-slice product and is not replaced by a statistical pilot error.

We now bound the pilots. Let \(S_0=s\wedge1\), let \(q_j=P_X(B_j)\), and let \(N\) denote the size of the pilot fold. Conditional on \(X\in B_j\), positivity and the interior choice of \(b\) give
\[
2bc\le P(|A-t_0|\le b\mid X\in B_j)\le 2b/c.
\]
Within a cube, the two slice functions vary by at most \(C\delta^{S_0}\); within the treatment window, the two treatment-direction functions vary by at most \(Cb^{A_0}\) and \(Cb^{B_0}\), respectively. Thus the conditional biases of the raw cell estimators are bounded by
\[
C(\delta^{S_0}+b^{A_0})\quad\text{and}\quad C(\delta^{S_0}+b^{B_0}).
\]
For completeness, the occupancy calculation does not require a lower bound on \(p_X\). If \(R_j\) is the number of pilot observations falling in the relevant cell and treatment window, then \(R_j\) is binomial with success probability between \(2bcq_j\) and \(2bq_j/c\). The elementary identities
\[
P(R_j=0)\le e^{-2Nbcq_j},\qquad
\mathbb E\left[\frac{\mathbf1\{R_j>0\}}{R_j^r}\right]\le \frac{C_r}{(Nbq_j)^r}\quad(r=1,2)
\]
when \(Nbq_j\ge1\), together with the trivial bound when \(Nbq_j<1\), imply after multiplication by \(q_j\) and summation over at most \(C\delta^{-d}\) cubes that the integrated second and fourth centered moments are bounded by \(C(Nb\delta^d)^{-1}\) and \(C(Nb\delta^d)^{-2}\), respectively. One way to verify the inverse-moment inequality is to split at half the binomial mean: below that point a Chernoff bound applies, and above it \(R_j^{-r}\le 2^r(Nbq_j)^{-r}\). The same calculation conditional on \(N_j\), followed by the binomial inverse-moment bound for \(N_j\), applies to the histogram treatment-density pilot. Clipping cannot increase its error because the true slice lies in \([c,1/c]\subset[c/2,2/c]\). Since the outcome and all pilots are uniformly bounded, the bias and centered-moment bounds combine to give
\[
(\mathbb E_PU^4)^{1/4}\le C\left\{\delta^{S_0}+b^{A_0}+(nb\delta^d)^{-1/2}\right\}=Cu,
\]
\[
(\mathbb E_PV^4)^{1/4}\le C\left\{\delta^{S_0}+b^{B_0}+(nb\delta^d)^{-1/2}\right\}=Cv.
\]
The constants are unchanged when fold sizes, rather than \(n\), appear because every fold size is proportional to \(n\).

It remains to pass from the population pilots to the empirical marginalizations and empirical local-polynomial fit. Conditional on the pilot fold, write
\[
\eta_m=\widehat m-P_X\widetilde\mu,\qquad \eta_f=\widehat f-P_X\widetilde\pi.
\]
The pilots are constant in the dose coordinate, so these quantities do not depend on \(a\). Boundedness and independence of the marginalization fold imply
\[
\mathbb E(\eta_m^2+\eta_f^2\mid I_\mu)\le C/n.
\]
Conditional on both training folds, the conditional mean \(g_{\mathcal F}(a)=\mathbb E[\widehat\xi\mid A=a,\mathcal F]\) obeys the exact decomposition
\[
g_{\mathcal F}(a)-\theta(a)=B(a)+\eta_m-\frac{\eta_f}{f(a)}P_X\left[\frac{v_a\pi_a}{\widetilde\pi_a}\right].
\]
The last display, boundedness, and the preceding variance inequality show that empirical marginalization contributes at most \(C/n\) to mean squared error.

Let \(N_{\mathrm{lp}}=|I_{\mathrm{lp}}|\). The population Gram has minimum eigenvalue at least \(\lambda_0\). Each summand in the empirical Gram has operator norm at most \(C/h\), and its centered second moment has operator norm at most \(C/h\). Applying the exponential-moment proof of matrix Bernstein to this fixed \((p+1)\)-dimensional sum yields
\[
P\{\lambda_{\min}(\widehat D_h)<\lambda_*\}\le C\exp(-c_1nh),
\]
for every fixed \(\lambda_*<\lambda_0\). This conclusion can also be obtained entrywise from scalar Bernstein followed by the fixed-dimensional operator-norm bound. On the good-Gram event, \(\|\widehat D_h^{-1}\|\le\lambda_*^{-1}\), and empirical polynomial reproduction is exact:
\[
e_0^\top\widehat D_h^{-1}\frac1{N_{\mathrm{lp}}}\sum_{i\in I_{\mathrm{lp}}}K_h(A_i)S_h(A_i)q(A_i)=q(t_0)
\]
for every polynomial \(q\) of degree at most \(p\). Apply this identity with \(q=Q_h\). The Taylor remainder contributes at most \(Ch^\alpha\) in root mean square because
\[
\mathbb E\left[\left\{N_{\mathrm{lp}}^{-1}\sum_i|K_h(A_i)|\right\}^2\right]\le C\left(1+\frac1{nh}\right).
\]
Conditional on the treatment values and the training folds, the centered pseudo-outcome score has variance bounded by
\[
\frac{C}{N_{\mathrm{lp}}^2}\sum_{i\in I_{\mathrm{lp}}}K_h(A_i)^2\|S_h(A_i)\|^2.
\]
Its expectation is at most \(C/(nh)\). The same bound covers random-design fluctuation of the bounded conditional mean. Combining these facts with the uniform bound on \(B\), and using Cauchy--Schwarz with the fourth-moment pilot bounds, gives
\[
\mathbb E_P[(\widehat\theta_{h,b,\delta}-\theta(t_0))^2;\,\lambda_{\min}(\widehat D_h)\ge\lambda_*]
\le C\left[h^{2\alpha}+\{(u+h^{A_0})(v+h^{B_0})\}^2+\frac1{nh}+\frac1n\right].
\]
The final clipping is nonexpansive relative to \(\theta(t_0)\in[-M_Y,M_Y]\). On the bad-Gram event both the clipped estimate and the target are in this interval, so its risk contribution is at most \(4M_Y^2C\exp(-c_1nh)\). The estimator is Borel measurable: cell counts and averages are Borel functions, clipping is continuous, the minimum eigenvalue is continuous in the matrix entries, and matrix inversion is continuous on the good-Gram set; the construction assigns zero before clipping on its complement. Thus it belongs to \(\mathfrak T_n\). Taking the supremum over \(P\) and using this single estimator in the minimax infimum proves the displayed master inequality.